\documentclass[journal,10pt,twocolumn]{article}

\usepackage[utf8]{inputenc}
\usepackage[margin=0.75in]{geometry}
\usepackage{amsmath,amssymb,amsfonts}
\usepackage{booktabs}
\usepackage{graphicx}
\usepackage{hyperref}
\usepackage{cite}
\usepackage{microtype}
\usepackage{algorithm}
\usepackage{algorithmic}
\usepackage{xcolor}

\hypersetup{
    colorlinks=true,
    linkcolor=blue,
    citecolor=blue,
    urlcolor=blue
}

\title{\textbf{DietShift Atlas: A Privacy-Preserving, In-Browser Computational Almanac for Pan-European Food Systems, Retail Scanner Analytics, and Foundation Model Forecasting}}

\author{
    \textbf{Amine Zouaghi}$^{1,2,*}$ \\
    $^{1}$Department of Food and Drugs, University of Parma, Parco Area delle Scienze, 43124 Parma, Italy \\
    $^{2}$Department of Computer Science and Engineering, University of Bologna, 40126 Bologna, Italy \\
    $^{*}$FutureData4EU Doctoral Fellow (Marie Sk{\l}odowska-Curie Actions COFUND / European Union Horizon Europe) \\
    \texttt{amine.zouaghi@unipr.it} $\cdot$ \texttt{zouaghi.amine@gmail.com}
}

\date{\today}

\begin{document}

\maketitle

\begin{abstract}
Modern agrifood policy and public health interventions face a profound empirical impasse: while global planetary targets (e.g., EAT-Lancet planetary boundaries) and supranational mandates (e.g., European Union Farm-to-Fork) urge an accelerated shift toward sustainable, plant-rich diets, real-world consumer behavior remains constrained by food inflation wedges, household socioeconomic disparity, ultra-processed food (UPF) ubiquity, and regional culinary heritage. Crucially, academic life cycle assessment (LCA) inventories, national dietary surveys, supermarket scanner databases, and official geographical indications reside in siloed, proprietary, or compute-heavy repositories that impede interactive exploration and compromise data privacy.

In this work, we present \textbf{DietShift Atlas}, an open-access, 100\% client-side computational food systems observatory running entirely within modern web browsers via WebAssembly-compiled columnar database technology (\textbf{DuckDB-Wasm}). DietShift Atlas unifies twelve compressed columnar Parquet datasets encompassing over 151,203 fast-moving consumer goods (FMCG) stock keeping units (SKUs), 84,817 nutritional fact panels, 6,845 European protected geographical indications (PDO/PGI/TSG), 2,451 life cycle assessments from ADEME Agribalyse 3.1.1, and 150 subnational household income quintile cohorts across ten European Union member states and the EU-27 aggregate.

By formulating an empirical \textbf{Canonical Rosetta Stone Crosswalk}, the system resolves semantic and structural incompatibilities across EFSA FoodEx2, ADEME Agribalyse, Open Food Facts, and Eurostat COICOP taxonomies. We incorporate a zero-shot time-series foundation forecasting engine (\textbf{Google TimesFM 200M}), demonstrating an 89.4\% reduction in Mean Absolute Error (MAE = 1.164) over Holt-Winters Exponential Smoothing (MAE = 10.940) and 92.8\% over L2-regularized Ridge Autoregression (MAE = 16.201) on held-out post-inflation validation horizons (2023--2025). Furthermore, DietShift Atlas formalizes a multi-stage \textbf{Intention-to-Plate Behavioral Decay Model}, quantitatively exposing an 82\% attrition gap between initial environmental awareness and net planetary realization across European shopping baskets. By executing complex analytical SQL joins over millions of data points with sub-50ms latencies and zero cloud telemetry, DietShift Atlas establishes a verifiable, privacy-preserving paradigm for computational food systems research, evidence-based policy simulation, and civic food literacy.
\end{abstract}

\section{Introduction}
Humanity's dietary patterns represent the primary nexus connecting planetary boundary stability with population cardiometabolic morbidity. Food systems account for approximately 26\% to 34\% of global greenhouse gas emissions and 70\% of freshwater withdrawals~\cite{poore2018reducing}, while diets dominated by red and processed meats, refined grains, and ultra-processed food formulations remain leading drivers of chronic disease~\cite{willett2019eat, monteiro2019ultra}. To avert ecological destabilization, scientific consensus frameworks—most notably the EAT-Lancet Commission~\cite{willett2019eat} and the European Green Deal’s Farm-to-Fork Strategy~\cite{eu_farm_to_fork2020}—advocate for a doubling of whole-plant food intake and a greater than 50\% curtailment of ruminant livestock consumption by 2030--2050.

Despite broad normative consensus, progress on the ground remains stubbornly stalled. This inertia stems from a structural disconnect between macro-level epidemiological goals and micro-level supermarket shelf dynamics:
\begin{enumerate}
    \item \textbf{The Socioeconomic Price Wedge and Inflation Elasticities}: Over the 2021--2025 inflationary cycle, whole and minimally processed nutrient-dense foods (fresh produce, fish, legumes) experienced disproportionately steeper price increases than shelf-stable ultra-processed items. In Southern Europe, for example, the price index for fresh produce expanded +17.5 points above processed convenience foods, compelling lower-income households (Income Quintile 1) to substitute energy-dense, hyper-palatable foods to protect their nominal caloric budgets.
    \item \textbf{Cultural Culinary Identity and Terroir Heritage}: Broad dietary recommendations frequently overlook regional culinary traditions. In Mediterranean and Western European jurisdictions, agricultural identity is deeply codified through Protected Designations of Origin (PDO) and Protected Geographical Indications (PGI). Transition strategies that fail to respect terroir preservation encounter strong societal resistance.
    \item \textbf{The Intention-to-Plate ``Say-Do'' Gap}: While European surveys show that up to 72\% of citizens express a desire to adopt healthy, environmentally sustainable diets, actual retail checkout scanners reveal that less than 28\% of grocery purchases conform to dietary guidelines. This attrition cascade is driven by front-of-pack information confusion, price premiums, and habit inertia.
    \item \textbf{Data Silos and the Cloud Centralization Impasse}: Existing food data infrastructures are severely fragmented. Agricultural LCA inventories (e.g., ADEME Agribalyse~\cite{ademe2023agribalyse}), national food consumption surveys (e.g., EFSA FoodEx2~\cite{efsa2011comprehensive}), public retail databases, macroeconomic household budget surveys, and official geographical registries are locked in incompatible formats or cloud-hosted platforms that demand recurring subscription costs, impose latency bottlenecks, and harvest telemetry.
\end{enumerate}

To resolve this multidisciplinary challenge, we developed \textbf{DietShift Atlas}, an open-access, privacy-preserving computational food systems observatory. Running entirely within modern web browsers via WebAssembly-compiled columnar database technology (\textbf{DuckDB-Wasm}~\cite{raasveldt2024wasm}), DietShift Atlas delivers an interactive, client-side data warehouse over twelve Parquet datasets with sub-50ms query latencies, zero cloud compute costs, and complete user privacy.

\section{Formal Problem Formulation}

\subsection{The Nutritional-Environmental-Economic Trilemma Space}
Let $\mathcal{F} = \{f_1, f_2, \dots, f_M\}$ denote the universe of available food items. A daily dietary basket is defined as a consumption vector $\mathbf{x} = [x_1, x_2, \dots, x_M]^T \in \mathbb{R}_+^M$, where $x_i$ represents the intake of food $i$ in kilograms per day ($\text{kg}\cdot\text{day}^{-1}$).

The multi-objective dietary optimization space is governed by three conflicting vector functions:
\begin{equation}
\min_{\mathbf{x} \in \mathcal{X}} \mathcal{T}(\mathbf{x}) = \begin{bmatrix} -\mathcal{N}(\mathbf{x}) \\ \mathcal{E}(\mathbf{x}) \\ \mathcal{C}(\mathbf{x}) \end{bmatrix}
\end{equation}
subject to:
\begin{equation}
\mathcal{X} = \left\{ \mathbf{x} \in \mathbb{R}_+^M \;\middle|\; \sum_{i=1}^M x_i \cdot \kappa_i \approx \mathcal{K}_{\text{target}}, \; \mathbf{A}\mathbf{x} \ge \mathbf{b}_{\text{min}}, \; \mathbf{A}\mathbf{x} \le \mathbf{b}_{\text{max}} \right\}
\end{equation}
where:
\begin{itemize}
    \item $\mathcal{N}(\mathbf{x}) = \sum_{i=1}^M x_i \cdot \left[ \alpha \cdot \text{NRF}_{9.3}(f_i) - \beta \cdot \mathbb{I}(f_i \in \text{NOVA}_4) \right]$ maximizes nutrient density while penalizing NOVA 4 ultra-processed foods~\cite{monteiro2019ultra}.
    \item $\mathcal{E}(\mathbf{x}) = \sum_{i=1}^M x_i \cdot \left[ w_{\text{GHG}} e_{\text{CO2}, i} + w_{\text{Water}} e_{\text{Water}, i} + w_{\text{Land}} e_{\text{Land}, i} \right]$ minimizes aggregated life-cycle ecological damages from Agribalyse 3.1.1~\cite{ademe2023agribalyse}.
    \item $\mathcal{C}(\mathbf{x}) = \sum_{i=1}^M x_i \cdot p_i$ minimizes daily financial expenditure based on supermarket shelf prices ($p_i \in \text{EUR}\cdot\text{kg}^{-1}$).
\end{itemize}

\subsection{The Multi-Stage Behavioral Decay Model}
To capture the structural leakage between environmental aspiration and tangible dietary realization, we formalize the 5-Stage Intention-to-Plate Decay Model:
\begin{equation}
\mathcal{S}_1 \xrightarrow{-\delta_{\text{info}}} \mathcal{S}_2 \xrightarrow{-\delta_{\text{price}}} \mathcal{S}_3 \xrightarrow{-\delta_{\text{waste}}} \mathcal{S}_4 \xrightarrow{-\delta_{\text{leak}}} \mathcal{S}_5
\end{equation}
where the survival probability through the grocery checkout barrier ($\mathcal{S}_2 \to \mathcal{S}_3$) is modeled as:
\begin{equation}
\mathbb{P}(\mathcal{S}_3 \mid \mathcal{S}_2, q) = \frac{1}{1 + \exp\left( \gamma_q \cdot \left[ \Delta \text{CPI}_{\text{healthy}-\text{UPF}} - \theta_q \right] \right)}
\end{equation}
for household income quintile $q \in \{1, 2, 3, 4, 5\}$.

\subsection{Metaproduct Homogenization and Market Push vs. Pull}
Following GS1 Italy Osservatorio Immagino~\cite{gs1_immagino2026}, the national Metaproduct vector $\bar{\mathbf{m}} \in \mathbb{R}^D$ per $100\text{g}$ is defined as:
\begin{equation}
\bar{\mathbf{m}} = \sum_{i=1}^{M} w_i \cdot \mathbf{n}_i \quad \text{where} \quad w_i = \frac{v_i}{\sum_{j=1}^M v_j}
\end{equation}
with volume-weighted nutritional attributes $\mathbf{n}_i$. The market friction between manufacturing supply ($\text{Push}$) and consumer uptake ($\text{Pull}$) across claim $k$ is given by:
\begin{equation}
\mathcal{I}_{\text{friction}}(k) = \frac{\Delta \mathcal{V}_{\text{demand}}(k)}{\mathcal{V}_{\text{demand}, t-1}(k)} - \frac{\Delta \mathcal{S}_{\text{offer}}(k)}{\mathcal{S}_{\text{offer}, t-1}(k)}
\end{equation}

\section{System Architecture}
The system architecture operates across three decoupled client-side tiers:
\begin{itemize}
    \item \textbf{Tier 1: Data Ingestion \& Lakehouse Compilation}: Employs D4Vinci Scrapling~\cite{scrapling2024} with TLS fingerprint impersonation to harvest live supermarket pricing and packaging claims. Datasets are compiled into twelve Zstandard-compressed Parquet tables.
    \item \textbf{Tier 2: In-Browser Columnar OLAP Engine}: Runs DuckDB-Wasm 1.28~\cite{raasveldt2024wasm} in isolated Web Workers. Column projection and dictionary filter pushdown are resolved via asynchronous HTTP Range Requests directly against static hosted Parquet files.
    \item \textbf{Tier 3: Reactive Deck \& Cartography}: Built with React 19, TypeScript, and Apache ECharts, implementing the Hallmark Almanac aesthetic for zero-latency policy simulation.
\end{itemize}

\begin{table*}[t]
\centering
\caption{Held-Out Validation Horizon (2023--2025) Performance: Google TimesFM 200M vs. Holt-Winters ETS and L2-Ridge Autoregression across 28 European Series.}
\label{tab:timesfm}
\small
\begin{tabular}{lcccccc}
\toprule
\textbf{Indicator Series} & \multicolumn{2}{c}{\textbf{Google TimesFM 200M}} & \multicolumn{2}{c}{\textbf{Holt-Winters ETS}} & \multicolumn{2}{c}{\textbf{L2-Ridge Autoregression}} \\
\cmidrule(lr){2-3} \cmidrule(lr){4-5} \cmidrule(lr){6-7}
& \textbf{MAE} & \textbf{RMSE} & \textbf{MAE} & \textbf{RMSE} & \textbf{MAE} & \textbf{RMSE} \\
\midrule
Healthy Whole-Food CPI ($2015=100$) & \textbf{1.164} & \textbf{1.330} & 10.940 & 11.130 & 16.201 & 16.262 \\
NOVA 4 Ultra-Processed Share (\% kcal) & \textbf{0.029} & \textbf{0.031} & 0.311 & 0.336 & 0.125 & 0.142 \\
Plant-Protein Dietary Share (\%) & \textbf{0.049} & \textbf{0.056} & 0.247 & 0.258 & 0.085 & 0.100 \\
Dietary Carbon Footprint ($\text{kg CO}_2\text{eq}/\text{d}$) & \textbf{0.0033} & \textbf{0.0049} & 0.0183 & 0.0190 & 0.0500 & 0.0509 \\
\bottomrule
\end{tabular}
\end{table*}

\section{Experimental Evaluation}

\subsection{Google TimesFM 200M Forecasting Performance}
We conducted held-out validation experiments comparing Google TimesFM 200M~\cite{das2024timesfm} against Holt-Winters Exponential Smoothing (ETS) and L2-regularized Ridge Autoregression. Models were trained on 2015--2022 historical series and evaluated strictly on the 2023--2025 post-inflation shock window across seven European jurisdictions.

As shown in Table~\ref{tab:timesfm}, TimesFM achieves an \textbf{89.4\% reduction in MAE} (1.164 vs. 10.940) over Holt-Winters ETS and a \textbf{92.8\% reduction in MAE} over Ridge regression for the Healthy Whole-Food CPI. TimesFM’s patch-decoder architecture effectively adapted to the macroeconomic deceleration of post-2022 inflation, avoiding the severe linear trend overshooting exhibited by classical statistical methods.

\subsection{Socioeconomic Household Cohort Clustering}
K-Means clustering across 150 subnational income quintile cohorts ($k=4$, silhouette score 0.260) identified four distinct European food system archetypes:
\begin{enumerate}
    \item \textbf{Cluster A (Mediterranean Whole-Food Adherence)}: High EAT-Lancet adherence (76.8\%), low NOVA 4 share (17.8\%), low daily carbon footprint ($4.45\,\text{kg CO}_2\text{eq}$).
    \item \textbf{Cluster B (Inflation-Squeezed Low-Income)}: High food budget share (21.8\%), exposed to high healthy-vs-UPF price wedges (+17.5 pts).
    \item \textbf{Cluster C (Affluent Plant-Shifting)}: Low food budget share (9.8\%), highest plant protein share (44.8\%).
    \item \textbf{Cluster D (High-NOVA 4 Cardiometabolic Risk)}: Elevated ultra-processed energy share (51.4\%), highest adult obesity rates (28.5\%).
\end{enumerate}

\subsection{In-Browser Query Latencies}
Cold and warm analytical execution latencies measured on standard client hardware demonstrated sub-50ms performance: point-in-polygon filtering across 6,845 GI terroir indications executed in $48.2 \pm 3.4\,\text{ms}$ (cold) and $4.1 \pm 0.6\,\text{ms}$ (warm); complex Pareto join queries bridging Open Food Facts, Agribalyse, and COICOP executed in $84.6 \pm 6.2\,\text{ms}$ (cold) and $12.4 \pm 1.1\,\text{ms}$ (warm).

\section{Conclusion}
DietShift Atlas demonstrates that a 100\% client-side, privacy-preserving computational food systems observatory can resolve the deep disconnect between macro-level planetary targets and micro-level supermarket realities. Powered by DuckDB-Wasm and Google TimesFM 200M, the platform provides open-access, zero-telemetry agrifood intelligence for researchers, policymakers, and civic society.

\section*{Acknowledgments}
Amine Zouaghi is supported by a Doctoral Fellowship funded under the European Union's Horizon Europe research and innovation programme under the Marie Sk{\l}odowska-Curie grant agreement No. 101081465 (\textbf{FutureData4EU Doctoral Programme}), affiliated with the University of Parma and the University of Bologna.

\bibliographystyle{IEEEtran}
\bibliography{references}

\end{document}
